\uuid{0e0x}
\chapitre{Statistique}
\niveau{L2}
\module{Analyse de données}
\sousChapitre{Estimation}
\titre{Estimateurs et intervalles de confiance}
\theme{statistiques}
\auteur{Maxime NGUYEN}
\datecreate{2022-09-07}
\organisation{AMSCC}
\difficulte{}
\contenu{
\texte{$\mathcal N(\mu,\sigma)$ désigne la loi normale de moyenne $\mu$ et d’écart-type $\sigma$. Les quantiles sont à gauche : $t_{\nu;\gamma}$ pour la loi $St(\nu)$, avec une probabilité $\gamma$ à gauche du quantile.}
\texte{On relève le taux de cholestérol de 20 personnes choisies indépendamment. On adopte un modèle normal $\mathcal N(\mu,\sigma)$.
\begin{center}\begin{tabular}{lrrrrrr}\hline
Taux (cg/L) & 120 & 160 & 200 & 240 & 280 & 320 \\
Effectif & 2 & 4 & 5 & 4 & 3 & 2 \\ \hline
\end{tabular}\end{center}}
\begin{enumerate}
\item \question{Estimer l'espérance et la variance. Préciser les estimateurs et leurs propriétés.}
\reponse{Les estimateurs $\overline X$ et $S^2=\sum(X_i-\overline X)^2/19$ sont sans biais et convergents dans ce modèle. Les réalisations sont $\bar x=216$ cg/L et $s^2=3604{,}21053$ (cg/L)$^2$, donc $s\simeq60{,}03508$ cg/L.}
\item \question{Construire un intervalle de confiance bilatéral à $99\%$ pour $\mu$.}
\indication{La variance est inconnue et le modèle normal est une hypothèse de l'énoncé : utiliser Student à 19 degrés de liberté.}
\reponse{Le pivot $(\overline X-\mu)/(S/\sqrt{20})$ suit exactement $St(19)$. Ainsi
\[\left[\bar x-t_{19;0{,}995}\frac s{\sqrt{20}};\ \bar x+t_{19;0{,}995}\frac s{\sqrt{20}}\right]
\simeq[177{,}594;254{,}406]\ \text{cg/L}.\]
La normalité ne découle pas du faible effectif ; elle est supposée ici.}
\end{enumerate}
}
